(1) 谐振子频率为 $\omega$ 时基态波函数 $\psi_0^{(\omega)}(x)=(m\omega/(\pi\hbar))^{1/4}\exp[-m\omega x^2/(2\hbar)]$。频率在一个理想瞬间改变只改变哈密顿量，不施加瞬时平移算符，故 $\psi(x,0^+)=\psi_0^{(\omega_i)}(x)$，通常并非新哈密顿量的基态。
(2) 初态二阶矩为 $\langle x^2\rangle=\hbar/(2m\omega_i)$、$\langle p^2\rangle=m\hbar\omega_i/2$。对新哈密顿量 $H_f=p^2/(2m)+m\omega_f^2x^2/2$，得 $\langle H_f\rangle=\hbar[\omega_i+\omega_f^2/\omega_i]/4$。由 $H_f=\hbar\omega_f(n_f+1/2)$，可解 $\bar n_f=\frac14(\omega_f/\omega_i+\omega_i/\omega_f-2)\ge0$。
(3) 两实高斯基态的重叠振幅为 $\langle0_f|0_i\rangle=\left[2\sqrt{\omega_i\omega_f}/(\omega_i+\omega_f)\right]^{1/2}$，所以 $P_0=2\sqrt{\omega_i\omega_f}/(\omega_i+\omega_f)$。原始基态是偶函数，新频率谐振子势仍具空间反演对称，奇宇称本征态与其内积严格为零，所以激发只能进入 $n_f=0,2,4,\ldots$。
(4) 在新哈密顿量下 Heisenberg 运动 $x(t)=x(0)\cos\omega_ft+p(0)\sin\omega_ft/(m\omega_f)$。初态 $\langle xp+px\rangle=0$，故 $\langle x^2(t)\rangle=\frac{\hbar}{2m}[\cos^2(\omega_ft)/\omega_i+(\omega_i/\omega_f^2)\sin^2(\omega_ft)]$。其振幅含 $\cos(2\omega_ft)$ 呼吸；若 $\omega_f=\omega_i$ 则方差恒定、$\bar n_f=0$，而任意 $\omega_f$ 时 $\langle H_f\rangle$ 在 $t>0$ 均守恒。
